% In this file you should put the actual content of the blueprint.
% It will be used both by the web and the print version.
% It should *not* include the \begin{document}

% This makes the macros `\inputleanmodule` and `\inputleannode` available.
\input{../../.lake/build/blueprint/library/BV}

\chapter{Bombieri--Vinogradov}

\section{Statement}

The main result is the Bombieri--Vinogradov theorem, which gives a
strong estimate for primes in arithmetic progressions on average over moduli.

\inputleannode{Bombieri-Vinogradov}
% \inputleannode{BombieriVinogradov.BV_Delta_1P}

\section{Assumptions}

We take Siegel--Walfisz and the large sieve inequality as axioms.

\inputleannode{siegel_walfisz}
\inputleannode{large_sieve}

\section{Definitions}

% \inputleannode{logIntegral}
\inputleannode{Delta}
\inputleannode{Delta_eq_sum_char}
\inputleannode{onCoprime}

\section{Reduction to \texorpdfstring{$\Delta_\Lambda$}{Delta Lambda}}

Bombieri--Vinogradov for $\psi$ is reduced to a bound for $\Delta_\Lambda$
via the following steps.

\inputleannode{BV_Delta_Lambda}

\inputleannode{Delta_Lambda_eq}
\inputleannode{sum_primes_not_dvd_log_eq_id}


% \inputleannode{BombieriVinogradov.Delta_primeIndicator_eq}
% \inputleannode{BombieriVinogradov.sum_primes_not_dvd_eq_li}
% \inputleannode{BombieriVinogradov.max_Delta_1P_le_max_Delta_Lambda}
% \inputleannode{BombieriVinogradov.BV_Delta_1P}
% \inputleannode{BombieriVinogradov.BV_Delta_Lambda}

\section{Vaughan's Identity}

The von Mangoldt function decomposes as $\Lambda = \Lambda^\sharp + \Lambda^\flat + \Lambda_{\le U}$,
where $UV \le \sqrt{x}$ and $U, V \ge e^{\sqrt{\log x}}$.

\inputleannode{LambdaSharp}
\inputleannode{LambdaFlat}
\inputleannode{LambdaLEU}
\inputleannode{Lambda_decomp}

\section{Small Primes (\texorpdfstring{$\Lambda_{\le U}$}{Lambda le U})}

The contribution of the small part $\Lambda_{\le U}$ is bounded trivially.

\inputleannode{sum_LambdaLEU_le}
\inputleannode{Delta_LambdaLEU_bound}
\inputleannode{BV_LambdaLE}

\section{Type I Sums (\texorpdfstring{$\Lambda^\sharp$}{Lambda sharp})}

The sharp part $\Lambda^\sharp$ is a type I sum, bounded via Abel summation.

\inputleannode{Delta_convolution_eq}
\inputleannode{Delta_one_bound}
\inputleannode{Delta_abel_summation}
\inputleannode{Delta_monotone_bound}
\inputleannode{Delta_flog_bound}
\inputleannode{Delta_LambdaSharp_bound}
\inputleannode{BV_LambdaSharp}

\section{Type II Sums (\texorpdfstring{$\Lambda^\flat$}{Lambda flat})}

\subsection{Definitions}

\inputleannode{S}
\inputleannode{T}

\subsection{Reduction to character sums}

\inputleannode{character_sum_by_conductor}
\inputleannode{character_sum_Mobius}
\inputleannode{Delta_LambdaFlat_decomp}
\inputleannode{Delta_LambdaFlat_small_conductor}
\inputleannode{BV_LambdaFlat_via_T}

\subsection{Large sieve estimates}

\inputleannode{LargeSieve_convolution}
\inputleannode{LambdaFlat_dyadic}
\inputleannode{BV_char_sum_bound}
\inputleannode{T_r_bound}
\inputleannode{BV_LambdaFlat}